\begin{helper}[Situation 7.3: a Gamma owner with a late actual witness]
Let $G$ be a finite loopless multigraph with a finite individual edge
type, let $[k]$ be its palette, and fix a distinguished edge $e_0$, an
input coloring, an exact ordered tree-sequence $T$, a generating coloring
$\varphi$, an ETT $\mathcal E$ on these data, a
\noderef{HierarchyBase} $\mathcal B$, a
\noderef{HierarchyCounterexample} $C$, and a
\noderef{HierarchyMinimumCounterexampleContext} $\mathit{ctx}$.
Thus the context supplies properness, exactly-one-uncoloredness, all six
technical assertions at strictly lower positive rungs, fewer-divider
induction at the same rung, rung-local lexicographic minimality, and the
terminal-ladder bound $|\overline\varphi(T_n)|\ge 2n+11$.
Write $n=\mathcal E.\mathit{rounds}$, $H=C.\mathit{hierarchy}$,
$q=H.\mathit{dividerCount}$, and set
\[
 A=\noderef{HierarchyStarredStage}(H,q),\quad
 \ell=\noderef{HierarchyFullTailLength}(C),\quad
 T(y_i)=\noderef{HierarchyFullTailPrefix}(C,i),\quad
 B_{\rm pen}=T(y_{\ell-1}),\quad y_\ell=T.\mathit{vertex}(\mathit{last}).
\]
Let $T_n$ be the terminal ladder tree
$\mathcal E.\mathit{series.data.tree}((\operatorname{Fin.last}n).\mathit{castSucc})$,
and put
\[
 D=\noderef{HierarchyStoredColorsAt}(H,q),\quad
 \Gamma=\bigcup_{\rho\in D}C.\mathit{good}.\Gamma^q_\rho,\quad
 M_0=\bigcup_{v\in V(\mathcal B.\mathit{core})\setminus V(T_n)}
             \noderef{MissingColors}(G,\varphi,v).
\]
Here $V(U)=\noderef{TreeSequenceVertexSet}(U)$ and
$\overline\varphi(U)=\noderef{TreeMissingColors}(\varphi,U)$.
Assume $2\le j<\ell$, where
$j=\noderef{HierarchyPathNumber}(H)$.
Let $\alpha$ be missing on $A$ and at $y_\ell$.
Suppose $\eta\in D$ and $\alpha\in C.\mathit{good}.\Gamma^q_\eta$.
Let $t$ be a natural number satisfying $1\le t<\ell$ and $j-1\le t$,
such that the totalized first $\eta$-missing position on $T$ has
index $|A|+t$ and $\eta$ is actually missing at
$y_t=T(y_t).\mathit{vertex}(\mathit{last})$ under $\varphi$.
Then there is an original-relative
\noderef{HierarchyReplacementPresentation} $P$ under whose coloring
$\eta$ is missing at both $y_t$ and $y_\ell$.
\end{helper}
\begin{proof}
Apply \noderef{HierarchyCounterexamplePenultimateElementary} to
$\mathit{ctx.proper}$, $\mathit{ctx.oneEdgeUncolored}$,
$\mathit{ctx.lowerRung}$, $\mathit{ctx.fewerDividers}$,
$\mathit{ctx.lexicographicMinimality}$, and
$\mathit{ctx.largeMissingSet}$, with the ambient data in the statement.
It makes $B_{\rm pen}$ elementary under $\varphi$.
The fields of \noderef{TashkinovHierarchy} and the entrywise definition
of \noderef{TreeSequenceInitialSegment} give $A\preccurlyeq T$.
Since $\ell=|T|-|A|\ge3$, the exact prefix $B_{\rm pen}$ has length
$|T|-1$ and contains $A$. The base core is below $A$: at $q=0$
it equals $A$, and at $q>0$ use the core-inside-first field followed
by the nested stages. Also $T_n\preccurlyeq\mathcal B.\mathit{core}$:
in PE this is the left-segment field of the ordered union in
\noderef{HierarchyBase}, and in non-PE the core equals $T_n$.
Consequently $T_n$ and the base core both lie in $B_{\rm pen}$.

By \noderef{GoodHierarchyReservedDisjoint}, every last Gamma pair lies
in $\overline\varphi(A)$ and $\Gamma\cap D=\varnothing$.
Expanding \noderef{HierarchyStoredColorsAt}, one also has
$D=D_n\setminus\overline\varphi(A)$, where
$D_n=\noderef{TashkinovStoredColors}(\mathcal E.\mathit{series.data},n)$.
In particular no stored color is missing on $A$.

Put $X=T(y_t)$. The bounds on $t$ give
$A\prec X\preccurlyeq B_{\rm pen}$, with $|X|=|A|+t$.
Since $\eta\in D$, it is not missing on $A$, whereas $\alpha$ is
missing there, so $\alpha\ne\eta$. The first $\alpha$ witness is on
$A$ and precedes the first $\eta$ witness $y_t$.
The avoidance clause of \noderef{GoodHierarchy} for
$\alpha\in C.\mathit{good}.\Gamma^q_\eta$, with
\noderef{HierarchySegmentColors}, excludes $\alpha$ at every exact
edge position $|A|\le r<|X|$.

Apply \noderef{ChangeLemma} to the ordered pair $(\alpha,\eta)$, with
the six context projections and the derived elementary prefix.
Both colors are missing on $B_{\rm pen}$. The conditional ordered-unused
premise is precisely the witness order and edge exclusion just proved.
The avoidance clause puts $\alpha$ in
\noderef{HierarchyStageMissingColors}(H,q)$: at positive divider count
this means missingness on $A=T_{n,q}$, and at zero it means missingness
on $T_n$. Thus the respective first or second conclusion gives equality
of the components through the first missing witnesses.
Apply \noderef{P2905_T2_ElementaryTerminalComponent} to those witnesses
and $w=y_\ell$, which misses $\alpha$. Both witnesses are in the
strict elementary prefix $B_{\rm pen}$ and different from $y_\ell$.
It gives the disjoint terminal component, its possibly singleton path
certificate, and pointwise equality of all prefix missing sets after
the swap
$\sigma=\noderef{TwoColorComponentSwap}(G,\varphi,\alpha,\eta,y_\ell)$.
The exchange clause of
\noderef{TwoColorComponentSwapMissingBehavior} makes $\eta$ missing at
$y_\ell$ and the prefix equality retains its actual witness $y_t$.
It also retains $\alpha$'s witness on $A$. Both affected colors are
therefore missing on $X$ under $\sigma$.

Apply \noderef{StableChangeCasesOneTwo} now to the same ordered colors,
context data, elementary prefix, missing witnesses, terminal path, swap,
Cases 1/2 alternative, and component disjointness. Retain its
\noderef{HierarchyWeaklyStable} conclusion at $A$ and its
\noderef{ModColoring} conclusion. We next verify all additional
SplitTail hypotheses rather than treating these two conclusions as a
full reconstruction.

\textbf{Fixed displayed interval and the exact seed exception.}
The absence of $\alpha$ fixes every $\alpha$-affected displayed edge
from $A$ to $X$. Apply \noderef{HierarchyTailMissingColorExclusion}
with the old elementary prefix, $i=t$, and the actual missing witness
of $\eta$ at $y_t$. For any $\eta$-colored displayed edge in this same
interval it gives precisely the exceptional data
\[
 q=0,\quad\mathcal B.\mathit{core}=T_n,\quad r=|A|,\quad
 \Theta_n\in\{SE,RE\},\quad F_n=\{f_n\},\quad
 T(y_t).\mathit{edge}(r)=f_n,\quad
 \varphi(f_n)=\eta\in D_n.
\]
Here $A=\mathcal B.\mathit{core}=T_n$. Inspect the actual terminal
certificate $\mathcal E.\mathit{series.certified}(\operatorname{Fin.last}n)$
of \noderef{TashkinovSeries}. Its SE constructor's
\noderef{MaximumDefectiveChoice} makes $f_n$ incident with $T_n$;
its RE constructor's
\noderef{RevisitingExtensionCertificate} gives an inside endpoint
in $T_n$ for the same individual edge. The connecting-edge singleton
identifies these constructor edges with the $f_n$ above.
The zero-divider non-PE clause of
\noderef{HierarchyWeaklyStable} is
$\noderef{StableColoring}(T_n,D_n,\varphi,\sigma)$.
Because $f_n$ is incident with $T_n$ and its old color is the protected
$\eta\in D_n$, the incident-edge clause gives
$\sigma(f_n)=\varphi(f_n)$. Every other edge in $A$ to $X$ has neither
affected color. This verifies fixed edge values on that entire interval,
including the exceptional seed.

\textbf{Actual route through $X$ and suffix factorization.}
Again inspect that actual terminal certificate. Its initial constructor
is impossible since the rung is positive. Its PE constructor gives the
old closure route from $T_n$ to the certified final series tree; its
SE/RE constructor gives the exact singleton $F_n=\{f_n\}$, the literal
\noderef{TreeSequenceAdjoinEdge} seed $S=T_n+f_n$, and the old closure
route from $S$ to that final tree. The
\noderef{ExtendedTashkinovTree} fields put the displayed $T$ below
the final series tree and strictly above $T_n$. Consequently the PE
source lies below $T$, and the SE/RE seed of length $|T_n|+1$ lies
below $T$: both are exact prefixes of the final tree, and their length
is at most $|T|$. Restrict those actual closure routes with
\noderef{TashkinovAugmentationSequenceRestrict}. Equivalently, these
are exactly the two alternatives of the registered principal
\noderef{P2813_T1_TerminalRoute}, whose statement additionally retains
the seed's incidence and color-control data. Denote its source by $S$,
taking $S=T_n$ in PE. The source is an exact prefix of $T$ by
\noderef{TashkinovAugmentationSequenceInitialSegment}, applied to
that actual old route.

In PE, $S\preccurlyeq A\prec X$ by the base's ordered-union
left segment and the hierarchy nesting. In SE/RE at $q>0$,
\noderef{HierarchyNonparallelFirstEdge} places the same exact connecting
edge at position $|T_n|$ of the first ordinary hierarchy stage.
The adjoin-edge data, common-prefix equality, and its length
$|T_n|+1$ put the seed below that stage and hence below $A$.
At $q=0$, non-PE gives $A=T_n$ and $t\ge1$ gives
$|S|=|A|+1\le|X|$; the seed and $X$ are exact prefixes of $T$, so
$S\preccurlyeq X$. This retains the actual SE/RE seed even at $t=1$.

Restrict the old route $S\to T$ at $X$ with
\noderef{TashkinovAugmentationSequenceRestrict}. Separately apply
\noderef{TashkinovAugmentationSequenceSuffix} to
$S\preccurlyeq X\preccurlyeq T$ and that old route to obtain the actual
old suffix $X\to T$. Restriction alone is not used to produce a suffix.

We transport the restricted route $S\to X$ under $\sigma$.
At $q>0$, \noderef{HierarchyWeaklyStable} includes
\noderef{TreeColoringInvariant} on $A$, so every displayed edge with
index at least one and below $|A|$ is fixed. The index-zero edge is
$e_0$ by \noderef{TreeSequence} and stays uncolored under the literal
swap, so it too is fixed. At $q=0$, its base clause is the actual combined
stable relation on $T_n$ and the base closure in PE, and the actual
stable relation on $T_n$ in non-PE. Properness of $\sigma$ is given by
\noderef{TwoColorComponentSwapProper}; the literal swap preserves the
value $\mathsf{none}$ and never changes a colored value to
$\mathsf{none}$, so it remains exactly-one-uncolored.
Apply \noderef{P2813_T1_BaseDisplayedEdges} to these original and new
one-uncolored laws and the base stable relation. It fixes every exact
displayed edge of $\mathcal B.\mathit{core}=A$ at $q=0$,
including the PE union entries not displayed in $T_n$.
Together with the already verified interval $A$ to $X$, all edge
positions of $X$ at or after $|S|$ are fixed. In zero-divider SE/RE
the route starts at the retained seed, not at $T_n$; its seed color
was separately fixed above.

Let $X^-$ be the exact \noderef{TreeSequencePrefix} of length $|X|-1$.
It is a positive-length tree-sequence since $A$ has positive length
and $t\ge1$, and it lies in $B_{\rm pen}$.
The pointwise missing-set equality on $B_{\rm pen}$ holds on $X^-$.
Apply \noderef{TashkinovAugmentationSequenceColoringTransport} with
source $S$, predecessor $X^-$, endpoint $X$, the restricted old route,
the predecessor inclusion and length equality, the fixed displayed
values at every position at or after $|S|$, and this missing-set equality.
It gives the genuine $\sigma$-route
\[
 T_n\to X\quad\hbox{in PE},\qquad
 T_n+f_n\to X\quad\hbox{in SE/RE}.
\]
The old suffix $X\to T$ is also transportable by
\noderef{TashkinovAugmentationSequenceComponentSwap}, since both
affected colors are missing on $X$ under $\sigma$; no fixed values
above $X$ have been asserted or used.

\textbf{All stored-first and Gamma cutoff conditions.}
The entire missing set on $B_{\rm pen}$ is now fixed. Therefore every
totalized \noderef{HierarchyFirstMissingPosition} on $T$ is fixed,
in particular that of every $\rho\in D$. Indeed, if a color has a
witness before the last vertex, equality at every preceding vertex
preserves exactly its nonempty set of early candidates and its least
candidate. If it has no early witness, neither coloring has one,
and both totalized positions are the last index, whether or not a
terminal witness exists. This checks the fallback as well as the
actual-witness case.

For each $\rho\in D$, suppose its last Gamma pair contains an
affected color. Reserved disjointness excludes $\eta$ from every such
pair. If the pair contains $\alpha$, its intersection with
$\Gamma^q_\eta$ is nonempty, so the disjoint-pairs clause of
\noderef{GoodHierarchy} forces $\rho=\eta$.
The old first position of this owner equals $|A|+t=|X|$ by the
hypothesis. Thus the cutoff condition holds for every owner, not merely
for the displayed $\eta$.

Apply \noderef{StableChangeCasesOneTwoSplitTail} with split prefix $X$.
We have checked its old prefix missingness, terminal path and swap,
Cases 1/2 data, component disjointness, strict $A$ to $X$ inclusion,
$X$ to $B_{\rm pen}$ inclusion, fixed interval (including the seed
exception), actual new route through $X$, actual old suffix, new
missingness of both colors on $X$, pointwise penultimate missing sets,
every stored first position, and the universal Gamma cutoff.
Retain its full corresponding maximum ETT, rebuilt base, hierarchy,
goodness, and same-good witness.
We describe the identity presentation used below. For every slot of
the certified \noderef{TashkinovSeries}, properness and exact
one-uncoloredness follow by cases on its
\noderef{TashkinovSeriesStep}: the initial case uses the two coloring
clauses of its \noderef{KTriple}; SE uses
\noderef{MaximumDefectiveChoice}'s selected-coloring clauses; PE uses
the explicit current-coloring clauses; RE copies the previous coloring.
Choose the same input, series, and ETT for the reflexive
\noderef{CorrespondingExtendedTashkinovTree}; use reflexive rung,
presentation, ladder, stored-set, connecting-edge, and kind equalities,
the original generating-coloring equality, the just verified slotwise
coloring laws, and reflexive \noderef{StableColoring} at each positive
slot. Define $P_0$ by coloring $\varphi$, that certificate, maximum
$C.\mathit{maximum}$, base $\mathcal B$, counterexample $C$, reflexive
\noderef{SameGoodHierarchy}, and the supplied context $\mathit{ctx}$.
These are all seven fields of
\noderef{HierarchyReplacementPresentation}.
Retain the provider's full corresponding certificate $\mathcal C$,
its maximum witness, base $\mathcal B'$, hierarchy $K$, good witness
$g$, and $\noderef{SameGoodHierarchy}(C.\mathit{good},g)$.
Construct a \noderef{HierarchyCounterexample} $C'$ on $\mathcal B'$:
its positive-rung field is $C.\mathit{positive}$ transported by
$\mathcal C.\mathit{rungAgreement}$, its hierarchy and goodness are
$K,g$, and its maximum is the provider's maximum witness.
The two stated missing witnesses, one in the strict prefix
$B_{\rm pen}$ and one at $y_\ell$, are distinct by injectivity of
\noderef{TreeSequence}'s vertex display. Their common missing color
contradicts the disjointness clause of
\noderef{ColoredVertexSetElementary}, proving the non-elementarity
field. Apply \noderef{HierarchyReplacementPresentationTrans} to $P_0$,
$\sigma$, $\mathcal C$, $\mathcal B'$, $C'$, and the supplied same-good
witness. It gives an original-relative presentation $P$ with coloring
exactly $\sigma$; in particular its minimum context is supplied by the
principal, not assumed.


Use the retained actual $\eta$-missing witness $y_t$ and the new
terminal $\eta$-missing witness in the counterexample construction.
The resulting presentation has coloring exactly $\sigma$, so it
has both asserted $\eta$-missing witnesses.
\end{proof}
